\chapter{Uniformly discrete spaces and bounded geometry}
\label{chap:bounded-geometry}

Milestone M8. A space has \emph{bounded geometry} if, at each scale, its balls are uniformly
small. There are two settings. For metric spaces we follow L\"oh (Definition 5.6.11), who
counts the points of closed balls and pairs bounded geometry with uniform discreteness. For
coarse spaces we follow J.~Roe, \emph{Lectures on Coarse Geometry}, AMS, 2003 (\S3.1), who
measures a set at the scale of a controlled set $E$, the \emph{gauge}, by its largest
$E$-separated subsets. A set $T$ is \emph{$E$-separated} if no two distinct points of $T$ are
$E$-related (\texttt{SetRel.IsSeparated} in Mathlib), and the \emph{$E$-ball} at $x$ is
$E_x = \{y \mid (y, x) \in E\}$ (\texttt{SetRel.ball}).

Uniform discreteness is a small-scale condition that coarse structures cannot see; Roe defines
it for metric spaces only (Example 3.14). On the coarse side its role is played by the
condition that the diagonal is a gauge. Uniformly discrete metric spaces are Mathlib's
\texttt{DiscreteUniformity} (design decision D22). The coarse counterpart and its name are an
open decision (D24); the coarse nodes follow its proposal.

\section*{Metric spaces}

\begin{definition}[Uniformly discrete space, L\"oh Def.~5.6.11]
  \label{DiscreteUniformity}
  A pseudometric space $X$ is \emph{uniformly discrete} if there is $\varepsilon > 0$ with
  $d(x, y) \ge \varepsilon$ for all $x \ne y$. In Lean this is Mathlib's class
  \texttt{DiscreteUniformity X}, which says that the diagonal is an entourage of the
  uniformity; for a pseudometric space the two agree, since the sets
  $\{(x, y) \mid d(x, y) < \varepsilon\}$, $\varepsilon > 0$, form a basis of the uniformity.
  \plannedlean{Metric.discreteUniformity_iff}
\end{definition}

\begin{definition}[Bounded geometry, L\"oh Def.~5.6.11]
  \label{Metric.HasBoundedGeometry}
  A pseudometric space $X$ has \emph{bounded geometry} if for every $r$ there is $K \in \N$
  such that every closed ball of radius $r$ in $X$ has at most $K$ points.
\end{definition}

L\"oh calls a metric space that is uniformly discrete and has bounded geometry a \emph{UDBG
space}; the conjunction gets no name here. Bounded geometry alone does not make a space
uniformly discrete: $\{n\} \cup \{n + 1/n\}$, $n \ge 1$, in $\R$ is a counterexample.

\begin{lemma}[Word metrics, cf.\ L\"oh after Def.~5.6.11]
  \label{Group.Generators.WithWordMetric.hasBoundedGeometry}
  \uses{Group.Generators.WithWordMetric, DiscreteUniformity, Metric.HasBoundedGeometry}
  A group with the word metric of a generating family is uniformly discrete. If the family is
  finite, the word metric has bounded geometry.
  \plannedlean{Group.Generators.WithWordMetric.instDiscreteUniformity}
\end{lemma}
\begin{proof}
  \uses{Group.Generators.finite_setOf_wordLength_le}
  Distances are word lengths, which are natural numbers, and only the identity has length
  $0$; so distinct elements are at distance at least $1$. Left multiplication by $g$ is an
  isometry taking the closed ball of radius $r$ about $1$ onto the one about $g$, so all closed
  balls of radius $r$ have as many points as $\{g \mid |g|_S \le \lfloor r \rfloor\}$, which is
  finite for a finite family (Lemma~\ref{Group.Generators.finite_setOf_wordLength_le}).
\end{proof}

\begin{lemma}[Coarse bounds on uniformly discrete spaces, cf.\ L\"oh Prop.~5.2.9]
  \label{Metric.CoarseLipschitzWith.exists_lipschitzWith}
  \uses{DiscreteUniformity, Metric.CoarseLipschitzWith, Metric.CoarseAntilipschitzWith}
  Let $f \colon X \to Y$ be a map of pseudometric spaces.
  \begin{enumerate}
    \item If $X$ is uniformly discrete and $f$ is coarse Lipschitz, then $f$ is Lipschitz.
    \item If $Y$ is uniformly discrete and $f$ is injective and coarse antilipschitz, then $f$ is
          antilipschitz.
  \end{enumerate}
  In particular an injective quasi-isometric embedding between uniformly discrete spaces is a
  bilipschitz embedding.
  \plannedlean{Metric.CoarseAntilipschitzWith.exists_antilipschitzWith}
\end{lemma}
\begin{proof}
  Let $\varepsilon > 0$ bound the distances between distinct points from below
  (Definition~\ref{DiscreteUniformity}). (1) If $d(f x, f y) \le L\,d(x, y) + C$, then for $x \ne
  y$, $C \le (C/\varepsilon)\,d(x, y)$, so $d(f x, f y) \le (L + C/\varepsilon)\,d(x, y)$; for $x =
  y$ both sides are $0$. (2) If $d(x, y) \le K\,d(f x, f y) + C$, then for $x \ne y$, $f x \ne f y$,
  so $C \le (C/\varepsilon)\,d(f x, f y)$, and $d(x, y) \le (K + C/\varepsilon)\,d(f x, f y)$.
\end{proof}

With Lemma~\ref{Group.Generators.WithWordMetric.hasBoundedGeometry}, this gives L\"oh's
Proposition 5.2.9: a bijective quasi-isometry between finitely generated groups with the word
metrics of finite generating families is a bilipschitz equivalence. Injectivity suffices, and the
Lipschitz bound needs neither injectivity nor a uniformly discrete codomain.

\section*{Coarse spaces}

\paragraph{Roe's definition.} Let $X$ be a coarse space. For $E \subseteq X \times X$ and
$S \subseteq X$, the \emph{$E$-capacity} $\operatorname{cap}_E(S) \in \N \cup \{\infty\}$ is the
largest size of an $E$-separated subset of $S$ (Roe, Definition 3.1). A \emph{gauge} is a
symmetric $E$ containing the diagonal $\Delta$. A set $F \subseteq X \times X$ is \emph{uniform}
with respect to a gauge $E$ if $\operatorname{cap}_E(F[E_x])$ and
$\operatorname{cap}_E(F^{-1}[E_x])$ are bounded uniformly in $x$, where
$F[A] = \{y \mid \exists a \in A,\ (y, a) \in F\}$ (Definition 3.5). The uniform sets form a
coarse structure $\mathcal{U}_E$, the \emph{universal bounded geometry structure} based on $E$
(Proposition 3.7, Definition 3.8), and $X$ has \emph{bounded geometry} if its coarse structure
is contained in $\mathcal{U}_E$ for some controlled gauge $E$ (Definition 3.9).

The definitions below state the same notion with four simplifications, listed after them with
the reason for each and the argument that nothing is lost. None of them is forced by Lean:
Roe's definition could be formalized as it stands. The first avoids proving a theorem before
the definition can be stated; the other three remove a construction or a hypothesis that the
definition does not need.

\begin{definition}[Bounded geometry with a gauge, cf.\ Roe Def.~3.5 and~3.9]
  \label{CoarseSpace.HasBoundedGeometryWith}
  \uses{CoarseSpace}
  Let $X$ be a coarse space and $E \subseteq X \times X$. The space has \emph{bounded geometry
  with gauge $E$} if for every controlled set $F$ there is $K \in \N$ such that for every
  $x \in X$, every $E$-separated subset of the ball $F_x$ has at most $K$ points.
\end{definition}

\begin{definition}[Bounded geometry, Roe Def.~3.9]
  \label{CoarseSpace.HasBoundedGeometry}
  \uses{CoarseSpace.HasBoundedGeometryWith}
  A coarse space has \emph{bounded geometry} if it has bounded geometry with some controlled
  gauge $E$.
\end{definition}

\subsection*{Differences from Roe's definition}

\begin{enumerate}
  \item \emph{No universal bounded geometry structure.} Roe says that the coarse structure is
        contained in $\mathcal{U}_E$. Stated literally in Lean, this needs $\mathcal{U}_E$ as
        a coarse structure (a \texttt{CoarseSpace}), and a coarse structure must be closed
        under products: that is Roe's Proposition 3.7, whose proof needs entropy as well as
        capacity (Lemma 3.6, Proposition 3.2(d)). The definition would then depend on a
        theorem. Containment only says that every controlled set is uniform with respect to
        $E$, and that is what Definition~\ref{CoarseSpace.HasBoundedGeometryWith} states, as a
        property of the coarse structure of $X$; no second coarse structure is built.
  \item \emph{No capacity function.} Roe measures sets by $\operatorname{cap}_E$, a supremum
        in $\N \cup \{\infty\}$. Here the definition says directly that every $E$-separated
        subset has at most $K$ points. Reason: the definition only ever compares a capacity
        with a bound, and a capacity function would need its own API (values in $\N \cup
        \{\infty\}$, suprema over subsets) with no other use yet; Mathlib has the metric
        version, \texttt{Metric.packingNumber}, but none for relations. Nothing is lost:
        $\operatorname{cap}_E(S) \le K$ exactly when every $E$-separated subset of $S$ has at
        most $K$ points (an infinite separated set has finite separated subsets of every
        size). If a later node needs capacities (growth, Roe's Proposition 3.15), the
        function can be added and this definition restated through it.
  \item \emph{No conditions on the gauge.} Roe's gauge is symmetric and contains the
        diagonal; here $E$ is any controlled set. Reason: the conditions are hypotheses every
        user would have to carry and check, and they do not change which sets are separated.
        A set is $E$-separated when no two \emph{distinct} points of it are $E$-related in
        either order (\texttt{SetRel.IsSeparated}); so adding $E^{-1}$ or $\Delta$ to $E$ does
        not change the $E$-separated sets. Hence if $E$ satisfies the definition, so does the
        controlled gauge $E \cup E^{-1} \cup \Delta$, with the same constants; and a gauge in
        Roe's sense is in particular a controlled set. The one place where Roe uses that a
        gauge contains $\Delta$ is item 4.
  \item \emph{Balls $F_x$ in place of $F[E_x]$ and $F^{-1}[E_x]$.} Roe bounds separated
        subsets of two sets for each controlled $F$; here, of one ball $F_x$. Reason: one
        condition instead of two, and no image of a set under a relation. Nothing is lost
        because the definition quantifies over \emph{all} controlled sets, which are closed
        under inverses and products: $F[E_x]$ is the ball at $x$ of $F \circ E$, and
        $F^{-1}[E_x]$ that of $F^{-1} \circ E$, both controlled when $E$ and $F$ are. So Roe's
        two conditions for $F$ are this definition's condition for $F \circ E$ and
        $F^{-1} \circ E$. Conversely $F_x \subseteq F[E_x]$, as $x \in E_x$ for a gauge.
        This step is valid only for the condition on all controlled sets at once, not for the
        uniformity of a single set $F$ (Roe's Definition 3.5 itself), which is why that
        notion is not defined here.
\end{enumerate}

Two further differences are not choices of this chapter. Roe's coarse structure is a
collection of controlled sets; Lean's \texttt{CoarseSpace} (Definition~\ref{CoarseSpace})
stores the filter of their complements, an equivalent encoding. And ``bounded uniformly in
$x$'' is written as a bound $K \in \N$ on cardinalities taken in $\N \cup \{\infty\}$, so that
infinite separated sets violate it. The equivalence with Roe's definition argued in items 1--4
is not formalized: it would need Roe's capacity, gauges and $\mathcal{U}_E$, which are
not defined here.

\begin{definition}[Uniformly locally finite coarse space, cf.\ Roe Ex.~2.44]
  \label{CoarseSpace.IsUniformlyLocallyFinite}
  \uses{CoarseSpace.HasBoundedGeometryWith}
  A coarse space is \emph{uniformly locally finite} if it has bounded geometry with the
  diagonal as gauge: for every controlled set $F$ there is $K \in \N$ such that every ball
  $F_x$ has at most $K$ points.
\end{definition}

Every set is separated for the diagonal, which gives the second form of the definition. The
diagonal is controlled, so a uniformly locally finite coarse space has bounded geometry.

\subsection*{Differences from Roe's Example 2.44}

Roe has no name for this condition. In Example 2.44 he defines, on a set $X$, the coarse
structure whose controlled sets $E$ have sections $E_x$ and $E^x = E^{-1}[\{x\}]$ of uniformly
bounded size, and Definition 3.8 identifies it with $\mathcal{U}_\Delta$. Here, as in item 1,
that coarse structure is not built: the definition says that every controlled set of the given
coarse structure satisfies Roe's condition, which is containment in $\mathcal{U}_\Delta$.
Only the sections $F_x$ are bounded, not $F^x$: the sections $F^x$ are the balls of the
controlled set $F^{-1}$ (item 4). The name ``uniformly locally finite'' is a proposal (design
decision D24), after the name of the same condition for graphs and simplicial complexes.

\section*{Comparison on metric spaces}

A pseudometric space carries the coarse structure in which a set of pairs is controlled when
the distance is bounded on it
(Definition~\ref{PseudoEMetricSpace.toCoarseSpace}).

\begin{lemma}[cf.\ Roe Ex.~3.14]
  \label{Metric.hasBoundedGeometry_iff_isUniformlyLocallyFinite}
  \uses{Metric.HasBoundedGeometry, CoarseSpace.IsUniformlyLocallyFinite,
    PseudoEMetricSpace.toCoarseSpace}
  A pseudometric space has bounded geometry if and only if its metric coarse structure is
  uniformly locally finite.
\end{lemma}
\begin{proof}
  \uses{Metric.isControlled_iff}
  The set $\{(x, y) \mid d(x, y) \le r\}$ is controlled, and its ball at $x$ is the closed ball
  of radius $r$ about $x$. Every controlled set is contained in one of these sets, so its balls
  are contained in closed balls of a fixed radius.
\end{proof}

\begin{theorem}[cf.\ Roe Ex.~3.13]
  \label{CoarseSpace.hasBoundedGeometry_iff_packingNumber_le}
  \uses{CoarseSpace.HasBoundedGeometry, PseudoEMetricSpace.toCoarseSpace}
  The metric coarse structure of a pseudometric space $X$ has bounded geometry if and only if
  there is $\varepsilon \ge 0$ such that for every $r$ there is $K \in \N$ with the following
  property: for every $x \in X$, every subset of the closed ball of radius $r$ about $x$ whose
  points are pairwise at distance greater than $\varepsilon$ has at most $K$ points. (The
  largest size of such a subset is Mathlib's packing number \texttt{Metric.packingNumber}.)
\end{theorem}
\begin{proof}
  \uses{Metric.isControlled_iff}
  For $\varepsilon \ge 0$ write $E_\varepsilon = \{(x, y) \mid d(x, y) \le \varepsilon\}$; a set
  is $E_\varepsilon$-separated exactly when its points are pairwise at distance greater than
  $\varepsilon$ (\texttt{Metric.isSeparated\_iff\_setRelIsSeparated}).
  ($\Rightarrow$) A controlled gauge $E$ lies in some $E_\varepsilon$, so every
  $E_\varepsilon$-separated set is $E$-separated; apply the definition to $F = E_r$, whose
  balls are the closed balls of radius $r$.
  ($\Leftarrow$) Take $E_\varepsilon$ as the gauge. A controlled $F$ lies in some $E_r$, so its
  balls lie in closed balls of radius $r$, and the bound for $r$ applies.
\end{proof}

Combining the two results with the remark after
Definition~\ref{CoarseSpace.IsUniformlyLocallyFinite}: a pseudometric space of bounded geometry
in L\"oh's sense has bounded geometry in Roe's sense (take $\varepsilon = 0$). The converse
fails, since bounded geometry in Roe's sense is invariant under coarse equivalence and L\"oh's
is not: $\R$ is coarsely equivalent to $\Z$.
